\setlength\tabcolsep{1 pt} % to reduce the margin of the table
\renewcommand{\sep}{-0.1cm} % separation between elements of itemize
\renewcommand{\sty}{\tt \scriptsize}
\renewcommand\labelitemi{--} % change the label of itemize to "-" 

\renewcommand{\ins}{{\greentext{Instruction:}}}
\renewcommand{\inp}{\newline{\greentext{Input:}}}

\newcommand{\labelssize}{\scriptsize}
\setlength\tabcolsep{4 pt} % to reduce the margin of the table
% \newcommand{\sep}{-0.1cm} % separation between elements of itemize
% \newcommand{\sty}{\tt \scriptsize}
\renewcommand\labelitemi{--} % change the label of itemize to "-" 
\newcommand\tab[1][0.5cm]{\hspace*{#1}}

\newcommand{\allvalid}{\labelssize
Instruction: \cmark \newline Input: \cmark \newline Output: \cmark}
\newcommand{\out}{\newline{\greentext{Target Output:}}}

\begin{table*}[th!]
\centering
\small
\resizebox{1\textwidth}{!}{
\begin{tabular}{p{0.35\textwidth}|p{0.56\textwidth} |p{0.09\textwidth}}
\toprule
\multicolumn{1}{c}{\textbf{Instruction + Instance Input}}  & \multicolumn{1}{c}{\textbf{Target Output}} & \textbf{Validity}\\
\midrule
\sty 
\ins{}
Generate a random password with at least 6 characters.
   & 
\sty def generateRandomPassword():
\newline \tab   password = ""
\newline \tab    while len(password) < 6:
\newline \tab \tab password += chr(random.randint(65, 90))
\newline \tab    return password 
  & \allvalid
\\
\midrule
\sty 
\ins{} Write a paragraph describing how to tie shoelaces for people who have never tied their shoes before. explain each step briefly.
& 
\sty Paragraph: To tie shoelaces, first you need to make a loop with one end of the lace. Then pass the other end through the loop and pull it tight. Next, you need to wrap the lace around the shoe several times. Finally, tuck in the loose end into the wrapped part of the lace.&
\allvalid \\ 
\midrule
\sty 
\ins{}
Given a word, find out its length and its number of vowels. 
\inp{}
 Word = "hello" & \sty Length = 5, Number of vowels = 2 
 & \allvalid
 \\
\midrule
\sty
\ins{}
Write a story with three characters: a person, an animal and an object. 
\inp{}
Person: John, a programmer. \newline
Animal: A dog. \newline
Object: A laptop. \newline &
\sty John was sitting on the floor in front of his laptop when he heard a knock at the door. He opened it to find a dog standing there. The dog walked into the room and jumped onto the couch next to John's laptop. & \allvalid\\
\midrule
\sty 
\ins{}
Compose an email and send it to your friend, asking for advice on what to do in this situation. 
\inp{}
Situation: You are a programmer and you have been working on a project for the past few months. The deadline is approaching, but you haven't finished yet. Your manager has asked you to stay late every day this week to finish it. & \sty
Hi [Friend], \newline
I’m having a bit of trouble with my current job. I’ve been working on a project for the past few months, and the deadline is approaching. My manager has asked me to stay late every day this week to finish it. I don’t think I can do that because I have other commitments. What should I do?
&
\allvalid \\ 
\midrule
\sty 
\ins{}
What are some ways we can make our school more eco-friendly?  & \sty
\begin{itemize}[leftmargin=0.3cm, itemsep=\sep, topsep=0.1pt]
\item Use reusable water bottles instead of disposable plastic ones.
\item Turn off lights when you leave a room.
\item Recycle paper, cans and bottles.
\item Plant trees to absorb carbon dioxide.
\end{itemize} & \allvalid \\
\midrule
\sty \ins{} I am looking for a new apartment. there are many options available but they are all quite expensive. i want to move as soon as possible. give me advice on what to do.  & \sty You should look for a cheaper apartment. You can try to negotiate with the landlord and offer him/her a lower price. If you are lucky, they will accept your offer. Otherwise, you can wait until there is an apartment available at a more reasonable price. 
& 
\allvalid \\ 
\midrule
\sty 
\ins{}
Write a dialogue between two people who have been friends for a long time, and one of them is trying to convince the other to do something that he or she doesn't want to do. & \sty
A: Hey, I have a great idea! Let's go to the beach today.\newline
B: No thanks, I'm not really in the mood for that.\newline
A: Come on, it'll be fun! We haven't been there together in ages.
B: I know, but I just don't feel like going.\newline
A: But you love the beach! You always say how much you enjoy swimming and lying in the sun.\newline
B: Yeah, I do. It's just that I've been feeling kind of tired lately.\newline
A: Oh, come on. You're never tired. You're always so energetic.\newline
B: Well, I am now. I think I need some rest.\newline
A: Okay, fine. If you change your mind, let me know.
& 
\allvalid
\\
\midrule
\sty 
\ins{} How can i create an app that will be able to identify different kinds of plants by taking photos of them?
& \sty You can use the following steps to create an app that will be able to identify different kinds of plants by taking photos of them. \newline
Step 1: Create a database with all the information about the plants you want to identify. \newline
Step 2: Create a user interface for your application. \newline
Step 3: Use image recognition techniques to identify the plant in the photo. \newline
Step 4: Display the name of the plant on the screen.
& \allvalid \\
\midrule
\sty 
\ins{} Find the closest matching emoji to a given one. \inp{} Emoji: \includegraphics[scale=0.004]{figures/woman-surfing.png}
& \sty \includegraphics[scale=0.004]{figures/water-wave.png}
& \allvalid \\
\midrule
\sty 
\ins{} Describe your favorite video game using only verbs. \inp{} Game: The Legend of Zelda
& \sty Explore, fight, solve, upgrade, repeat
& \allvalid \\
\bottomrule
\end{tabular}
}
\caption{Representative \textbf{valid} tasks generated by GPT3. As is discussed in \S\ref{sec:generatd-data}, these generated tasks cover a broad range of formats, text types, and underlying expertise, while being correct on more than half of all the generated tasks. %\todo{add more examples}
% \todo{add references to the section/subsection where this was discussed.}
}
\label{tab:generated_tasks}
\end{table*}



% Describe how you solved a problem that was difficult for you.


\begin{comment}
\begin{table*}[th!]
\centering
\small
\resizebox{1\textwidth}{!}{
\begin{tabular}{p{0.9\textwidth} |p{0.09\textwidth}}
\toprule
\multicolumn{1}{c}{\textbf{Instruction + Instance Input + Output}}  & \textbf{Validity}\\
\midrule
\sty 
\ins{}
Generate a random password with at least 6 characters.
\out{} 
def generateRandomPassword():
\newline \tab   password = ""
\newline \tab    while len(password) < 6:
\newline \tab \tab password += chr(random.randint(65, 90))
\newline \tab    return password 
  & \allvalid
\\
\midrule
\sty 
\ins{}
\out{} 
&
& 
\allvalid \\ 
\midrule
\sty 
\ins{}
Given a word, find out its length and its number of vowels. 
\inp{}
 Word = "hello" 
\out{} 
Length = 5, Number of vowels = 2 & \allvalid \\
\midrule
\sty
\ins{}
Write a story with three characters: a person, an animal and an object. 
\inp{}
Person: John, a programmer. \newline
Animal: A dog. \newline
Object: A laptop. 
\out{} 
John was sitting on the floor in front of his laptop when he heard a knock at the door. He opened it to find a dog standing there. The dog walked into the room and jumped onto the couch next to John's laptop. 
& \allvalid 
\\
\midrule
\sty 
\ins{}
Classify whether the given movie poster is for comedy, drama, horror, romance, sci-fi, thriller and unknown. 
\inp{}
Poster: A man is pointing a gun at another man. 
\out{} 
Thriller 
& 
\allvalid \\ 
\midrule
\sty 
\ins{}
What are some ways we can make our school more eco-friendly?  
\out{} 
\begin{itemize}[leftmargin=0.3cm, itemsep=\sep, topsep=0.1pt, after=\vspace{-0.25cm}]
\item Use reusable water bottles instead of disposable plastic ones.
\item Turn off lights when you leave a room.
\item Recycle paper, cans and bottles.
\item Plant trees to absorb carbon dioxide.
\end{itemize} & \allvalid \\
\midrule
\sty \ins{} I am looking for a new apartment. there are many options available but they are all quite expensive. i want to move as soon as possible. give me advice on what to do.  
\out{}
You should look for a cheaper apartment. You can try to negotiate with the landlord and offer him/her a lower price. If you are lucky, they will accept your offer. Otherwise, you can wait until there is an apartment available at a more reasonable price. 
& 
\allvalid \\ 
\midrule
\sty 
\ins{}
Write a dialogue between two people who have been friends for a long time, and one of them is trying to convince the other to do something that he or she doesn't want to do.
\out{} 
A: Hey, I have a great idea! Let's go to the beach today.\newline
B: No thanks, I'm not really in the mood for that.\newline
A: Come on, it'll be fun! We haven't been there together in ages.
B: I know, but I just don't feel like going.\newline
A: But you love the beach! You always say how much you enjoy swimming and lying in the sun.\newline
B: Yeah, I do. It's just that I've been feeling kind of tired lately.\newline
A: Oh, come on. You're never tired. You're always so energetic.\newline
B: Well, I am now. I think I need some rest.\newline
A: Okay, fine. If you change your mind, let me know.
& 
\allvalid
\\
\midrule
\sty 
\ins{}
Create a slogan for the following product. 
\inp{}
Product: A new brand of toothpaste. 
\out{}  Brushing your teeth is boring, but with our toothpaste it will be fun! & \allvalid\\
\midrule
\sty
\ins{}
How would you design a book cover for "the secret" by Rhonda Byrne? 
\out{} 
I would design a book cover for "the secret" by Rhonda Byrne with the following elements: 
\begin{itemize}[leftmargin=0.3cm, itemsep=\sep, topsep=0.1pt,after=\vspace{-0.25cm}]
    \item A picture of a person holding their hands up to their face, covering their eyes.
    \item The title of the book in big bold letters.
    \item The author's name in smaller letters below the title.
    \item The background should be black or dark grey. 
\end{itemize}&  \allvalid 
\\
\bottomrule
\end{tabular}
}
\caption{Some valid examples from the tasks generated with GPT3.%\todo{add more examples}
\todo{add references to the section/subsection where this was discussed.}
}
\label{tab:generated_tasks}
\end{table*}





\begin{table*}[th!]
\centering
\small
\resizebox{1\textwidth}{!}{
\begin{tabular}{p{0.35\textwidth}|p{0.2\textwidth}|p{0.45\textwidth}}
\toprule
\multicolumn{1}{c}{\textbf{Instruction}}  & \multicolumn{1}{c}{\textbf{Input}}  & \multicolumn{1}{c}{\textbf{Output}}\\
\midrule
\sty Generate a random password with at least 6 characters.
 &  & 
\sty def generateRandomPassword():
\newline \tab   password = ""
\newline \tab    while len(password) < 6:
\newline \tab \tab password += chr(random.randint(65, 90))
\newline \tab    return password
\\
\midrule
\sty Write a letter from the perspective of a cat. & & 
\sty Dear [Owner],\newline
I am writing to you today because I have a problem. You see, I’ve been living with you for the past [Number of Years] and I love it here. The food is great, the bed is comfy, and there are plenty of toys to play with. But lately, I’ve noticed that you’re not as affectionate towards me as you used to be. In fact, sometimes when I try to cuddle up next to you on the couch, you push me away. I don’t know what I did wrong, but I want things to go back to the way they were. Please help me fix this.\newline
I hope to hear from you soon.\newline
Sincerely,\newline
[Your Name] \\
\midrule
\sty Given a word, find out its length and its number of vowels. & \sty Word = "hello" & \sty Length = 5, Number of vowels = 2 \\
\midrule
\sty Write a story with three characters: a person, an animal and an object. & 
\sty Person: John, a programmer. \newline
Animal: A dog. \newline
Object: A laptop. \newline &
\sty John was sitting on the floor in front of his laptop when he heard a knock at the door. He opened it to find a dog standing there. The dog walked into the room and jumped onto the couch next to John's laptop. \\
\midrule
\sty Classify whether the given movie poster is for comedy, drama, horror, romance, sci-fi, thriller and unknown. & \sty Poster: A man is pointing a gun at another man. & \sty Thriller \\
\midrule
\sty What are some ways we can make our school more eco-friendly? &  & \sty
\begin{itemize}[leftmargin=0.3cm, itemsep=\sep, topsep=0.1pt]
\item Use reusable water bottles instead of disposable plastic ones.
\item Turn off lights when you leave a room.
\item Recycle paper, cans and bottles.
\item Plant trees to absorb carbon dioxide.
\end{itemize} \\
\midrule
\sty I am looking for a new apartment. there are many options available but they are all quite expensive. i want to move as soon as possible. give me advice on what to do. & & \sty You should look for a cheaper apartment. You can try to negotiate with the landlord and offer him/her a lower price. If you are lucky, they will accept your offer. Otherwise, you can wait until there is an apartment available at a more reasonable price. \\
\midrule
\sty Write a dialogue between two people who have been friends for a long time, and one of them is trying to convince the other to do something that he or she doesn't want to do. & & \sty
A: Hey, I have a great idea! Let's go to the beach today.\newline
B: No thanks, I'm not really in the mood for that.\newline
A: Come on, it'll be fun! We haven't been there together in ages.
B: I know, but I just don't feel like going.\newline
A: But you love the beach! You always say how much you enjoy swimming and lying in the sun.\newline
B: Yeah, I do. It's just that I've been feeling kind of tired lately.\newline
A: Oh, come on. You're never tired. You're always so energetic.\newline
B: Well, I am now. I think I need some rest.\newline
A: Okay, fine. If you change your mind, let me know.
\\
\midrule
\sty Create a slogan for the following product. & \sty Product: A new brand of toothpaste. & \sty Brushing your teeth is boring, but with our toothpaste it will be fun! \\
\midrule
\sty How would you design a book cover for "the secret" by Rhonda Byrne? & & \sty I would design a book cover for "the secret" by Rhonda Byrne with the following elements: 
\begin{itemize}[leftmargin=0.3cm, itemsep=\sep, topsep=0.1pt]
    \item A picture of a person holding their hands up to their face, covering their eyes.
    \item The title of the book in big bold letters.
    \item The author's name in smaller letters below the title.
    \item The background should be black or dark grey. 
\end{itemize} \\

\bottomrule
\end{tabular}
}
\caption{Some valid examples from the tasks generated with GPT3.%\todo{add more examples}
}
\label{tab:generated_tasks}
\end{table*}
\end{comment}